\subsection{Notation}

Let \(R = \Fq[x_1, \dots, x_n]\) denote the polynomial ring in \(n\) variables over the finite field \(\Fq\). Given a polynomial \(f \in R\), we use \(\lm(f), \lc(f), \lt(f)\) to denote the leading monomial, the leading coefficient and the leading term (\(\lt(f) = \lc(f)\cdot \lm(f)\)) of the polynomial \(f\). We use $d_i$ to denote the maximum degree of the polynomials with respect to the target variable $x_i$. 

\subsection{Preliminaries}
Because we are exclusively interested in finding roots over the finite field $\Fq$, any valid solution must inherently satisfy the field equations $x_i^q - x_i = 0$ for all $i \in \{1, \dots, n\}$. Consequently, it is not restrictive to consider our polynomial systems modulo the ideal generated by these field equations. Throughout this paper, when manipulating polynomial equations, we implicitly reduce them with respect to $x_i^q - x_i$. This foundational reduction ensures that the degree of any polynomial with respect to a single variable $x_i$ is strictly bounded, such that $d_i \le q-1$. Establishing this bound early is critical, as it guarantees that the polynomial degrees remain tractable and prevents exponential degree growth during successive algebraic eliminations.